\scanpage{230}
% horizontal sections; three displays
sections horizontales au voisinage de $s_0$, telles que
\[
 a_0(s_0)=a(s_0),\quad b_0(s_0)=b(s_0).
\]
On a
\[
 \mu(a,b)=\underbrace{\mu(a_0,b_0)}_{
  \deleted{\int_{a_0}^{\unsure{b_0|b}}\varphi}\ \varphi(b_0)-\varphi(a_0)}
 +\mu(b_0,b)-\mu(a_0,a)
\]
donc
\[
 \mu(a,b)(s_0)=\bigl[\varphi(b(s_0))-\varphi(a(s_0))\bigr]
 +\left(\sum_i D_{i\,s}\int_{b_0(s)}^{b(s)}\omega_{i\,s}
 -\sum_i D_{i\,s}\int_{a_0(s)}^{a(s)}\omega_{i\,s}\right)_{s=s_0}
\]
Il faut calculer les deux termes de \unsure{rabiot}, c'est une histoire purement locale.

% lower half cancelled by large diagonal strokes
\deleted{Notons qu'on a, grâce à la connexion de $X/S$
\[
 da\in\Gamma(S,a^*(\unsure{\mathscr T|\mathcal T|\mathscr V}_{X/S})\otimes\Omega^1_S)
\]
d'où, en contractant avec
\[
 a^*(\omega)\in\Gamma(S,\underbrace{a^*(\Omega^1_{X/S})}_{
 \substack{\text{dual de}\\a^*(\mathscr T_{X/S})}})
\]
% slanted, outside the diagonal cancellation on the manuscript
}
\margin{donner une formule..... !}
\deleted{une section, que nous noterons $a^*(\omega)$
\[
 a^*(\omega)\in\Gamma(S,\Omega^1_S)
\]
Cette \deleted{section} \inserted{forme} est encore \uline{fermée}, \inserted{$\omega$ l'étant} [N.B. pour le voir définir $a^*(\omega)$ pour tte forme (degré quelconque, fermée ou non) sur $X$, et noter que $a^*$ commute à la différentielle ext.)}
% archive sheet number: 229
